\documentclass[xcolor=table]{beamer}
\usepackage[utf8]{inputenc}
\usepackage[T1]{fontenc} 
\usepackage[english]{babel} 
\usepackage{graphicx}
\usepackage{comment}
\usepackage{graphicx,wrapfig,lipsum}
\usepackage{hyperref}
\hypersetup{
    colorlinks=true,
    linkcolor=blue,
    filecolor=magenta,      
    urlcolor=cyan,
}
\setbeamertemplate{footline}[text line]{%
  \parbox{\linewidth}{\vspace*{-8pt}\today\hfill\insertshortauthor\hfill\insertpagenumber}}{}

%%Defining the ``proposition'' environment
\newtheorem{proposition}{Proposition} 

%%Sets the beamer theme "Cuerna"
\usetheme{Cuerna} 

\usecolortheme{default}
%Available color themes: default, bluesimplex, brick, lettuce

%%Insert the logo
\logo{\includegraphics[width=1.7cm]{plots/logo.jpeg}}

%%Title
\title{TELLIE PCA \\ Automation}

\author{Michal Rigan\\ %Author
          \texttt{mrigan@snolab.ca}} %e-mail
          
\date{\textbf{Overview}\\
\today} %Date or event

\institute{University of Sussex} %%Institution

\begin{document}

{
\setbeamertemplate{footline}{} 
\begin{frame}
  \titlepage %Creates the title page
\end{frame}
}

\begin{frame}{This study}
\begin{itemize}
	\item \textbf{Data-taking Auto}
	\begin{itemize}
		\item Profiling (how often? once before set?)
		\begin{itemize}
			\item Calibration curves (IPW - PIN)
			\item Tuning curves (PIN - NHit)
		\end{itemize}
		\item Data-recording
		\begin{itemize}
			\item Mode: \;\;\;\;\;\;\;\;\;\;\; a) dedicated \;\;\;\;\; | \;\;\;\;\; b) continuous
			\item Environment: \;\;\;\;\; a) ORCA \;\;\;\;\; | \;\;\;\;\; b) \underline{python}
		\end{itemize}
	\end{itemize}
	\item \textbf{Processing Auto}
\end{itemize}
\end{frame}

\begin{frame}{Processing Automation}
\textbf{GOAL:}\\script that reads run list for 95 fibres, then does the processing and analysis automatically, with monitoring.\newline \newline
This way it would work for both continuous and dedicated mode. The selection / creation of the list can also be automated (later).
\end{frame}

\begin{frame}{Steps}
\textbf{What data:}\\
(IF DEDICATED) Need a method to select the best run when we have retakes. Could use the nearline processors that we have (may need modifying) and fits.\newline \newline
(IF CONTINUOUS) Just take the latest run for each fibre. Need to create the list once we have collected data. \\

What to do if run type is changes - say ECA? Hold on to data-taking, then resume?
\end{frame}

\begin{frame}{Steps}
\textbf{Access to data:}\\
Options:
\begin{itemize}
	\item Move the reprocessing responsibility onto the processing group. Then need to download the data, which is almost always complicated and will be hard to automate.
	\item Do it ourselves: apollo2 OR use SNO+ computers. Data is retained UG for at least a week (is that enough?). SNUG2 is unused and available. Do processing ourselves.
\end{itemize}
\end{frame}

\begin{frame}{Steps}
\textbf{PCA table:}\\
PCA tables contains data and fits needed for PCA extraction.\\
Requires:
\begin{itemize}
	\item pca offsets = the injection time.\\Fit to residual hit times (hit - ToF - bucket).
	\item fitted position*
	\item fitted direction*
	\item angular systematic fits*
	\item * depends on pca offset values
\end{itemize}
Once calculated, create and update final json table.\\
\textbf{UNIFIED event selection, checks, cuts, ranges...}
\end{frame}

\begin{frame}{Steps}
\textbf{Extract constants:}\\
= running the PCA processor. Still running custom version, needs PR. Split LB - TELLIE proc?\\
Requires:
\begin{itemize}
	\item PCA table
	\item recent ECA
\end{itemize}
When/how to store and load the constants? Currently starting from next run once calib is finished - in continuous this can be a week late...\\
If we decide to load when the calib was done, will need a table for each fibre - 95 tables to load...
\end{frame}

\begin{frame}{Steps}
\textbf{Benchmarking:}\\
Maaany things to look at:
\begin{itemize}
	\item compare PCA table values
	\item apply PCA constants to a run, fit residual hit times, compare distribution to previous
	\item fit residual prompt peak
	\item compare TW
	\item compare cable delays
	\item other related fits: charges, LED offsets, RMS...
\end{itemize}
\end{frame}

\begin{frame}{Steps}
\textbf{Monitoring:}\\
\begin{itemize}
	\item Have a PCA page (has not been used or updated). It does what we need, is reeeeeeal slow. \newline \newline
	\item Will need logs / monitoring for the different steps of processing once automated - using couchdb, using some overview platform similar to BT.
\end{itemize}
\end{frame}

\end{document}